% Generated draft: requires mathematical and source review.
\documentclass[11pt]{article}
\usepackage[margin=25mm]{geometry}
\usepackage{fontspec}
\setmainfont{Latin Modern Roman}
\usepackage{amsmath,amssymb,mathtools,unicode-math}
\setmathfont{Latin Modern Math}
\usepackage{enumitem,needspace,xurl,hyperref}
\setlength{\parindent}{0pt}
\setlength{\parskip}{5pt}
\setlength{\emergencystretch}{3em}
\newcommand{\N}{\mathbb{N}}
\newcommand{\Z}{\mathbb{Z}}
\newcommand{\Q}{\mathbb{Q}}
\newcommand{\R}{\mathbb{R}}
\newcommand{\C}{\mathbb{C}}
\newcommand{\F}{\mathbb{F}}
\newcommand{\A}{\mathbb{A}}
\newcommand{\PP}{\mathbb{P}}
\newcommand{\E}{\mathbb{E}}
\newcommand{\eps}{\varepsilon}
\newcommand{\dd}{\,\mathrm{d}}
\begin{document}
\section*{UnsolvedMath ID: 30001504}
\textbf{Source problem code:} catalogue entry
\subsection{Definitions and mathematical statement}
For an integer \(d\geq 1\), let \(\mathbb T^d=\mathbb R^d/\mathbb Z^d\), with its quotient topology. An element \(A\in\operatorname{SL}_d(\mathbb Z)\) induces the toral automorphism \(x+\mathbb Z^d\mapsto Ax+\mathbb Z^d\). Call this automorphism totally irreducible if, for every integer \(n\geq 1\), the linear map \(A^n\) preserves no nonzero proper \(\mathbb Q\)-linear subspace of \(\mathbb Q^d\). For an abelian group \(H\), write \(\operatorname{rank}_{\mathbb Z}H=\dim_{\mathbb Q}(H\otimes_{\mathbb Z}\mathbb Q)\), regarding \(H\) as a \(\mathbb Z\)-module.

The supplied normalized formulation considers abelian subgroups \(G\leq\operatorname{SL}_d(\mathbb Z)\) with \(\operatorname{rank}_{\mathbb Z}G\geq 2\), containing a totally irreducible element, and maximal in rank among abelian subgroups containing them. Reading the last condition as maximality among supergroups of \(G\), it means
\[
\operatorname{rank}_{\mathbb Z}H\leq\operatorname{rank}_{\mathbb Z}G
\quad\text{for every abelian }H\text{ with }G\leq H\leq\operatorname{SL}_d(\mathbb Z).
\]
The diagonal action is
\[
g\cdot(x,y)=(gx,gy),\qquad g\in G,\quad (x,y)\in\mathbb T^d\times\mathbb T^d.
\]
\textbf{Classification problem.} For every such \(d\) and \(G\), classify all closed subsets of the form
\[
\overline{G\cdot(x,y)}
=\overline{\{(gx,gy):g\in G\}},
\qquad (x,y)\in\mathbb T^d\times\mathbb T^d,
\]
where closure is taken in the product topology.

\textbf{Unresolved source qualification.} The original wording asks this question for a ``full Cartan action'' by a group satisfying ``Condition 1.'' Neither Condition 1 nor the source definition of full Cartan action is supplied. Consequently, equivalence with the normalized hypotheses above cannot be established from this record. The phrase ``containing it'' in the normalized statement also leaves the referent of the maximal-rank condition uncertain: the display above uses supergroups of \(G\), whereas the wording could refer to abelian groups containing the totally irreducible element. This draft makes that reading explicit without resolving the source ambiguity.
\subsection{Short English statement}
Classify the closures of all diagonal orbits on the product of two identical tori under the stated abelian matrix actions. The original reference to Condition 1 and the precise maximal-rank hypothesis require source review. The supplied partially solved label and literature assessment are dated curation, not independently verified conclusions.
\subsection{Sources}
\begin{itemize}
\item Local UnsolvedMath registry, \texttt{data/unsolvedmath/problems.json}, permanent ID \(30001504\), title ``Orbit Closures of Full Cartan Torus Actions,'' exact source problem code \texttt{catalogue entry}. No individual proposer is supplied; the proposed year is recorded as 2010.
\item Supplied report citation: \emph{Homogeneous Dynamics and Number Theory (2011)}, Oberwolfach Reports, pp.\ 1703--1759, DOI \texttt{10.4171/owr/2010/29}. Source URL and literature-source label ``original OWR report'': \url{https://doi.org/10.4171/owr/2010/29}. The background describes these reports as records of mathematical workshops and their open-problem discussions.
\item Supplied literature source, labeled ``Lindenstrauss Wang'': \url{https://doi.org/10.1215/00127094-1593290}. No paper title or fuller bibliographic citation is supplied.
\item Literature assessment dated exactly \texttt{2026-08-21}: the supplied curation labels the problem partially solved and reports that Lindenstrauss and Wang classify topologically transitive self-joinings in rank at least three and construct a rank-two counterexample, leaving the unrestricted formulation partially resolved. This is supplied dated curation derived from the Oberwolfach material, not independently verified here; its status and cited claims require independent verification.
\item The registry labels the normalized statement \texttt{corrected\_verified} and reports that its mathematical content was checked against the official Oberwolfach report, with notation, context, and question form normalized without changing the source-level claim. This verification assertion is supplied metadata, not a verification performed for this draft.
\end{itemize}
\end{document}
